\section{Introduction}\label{sec:intro}

A US trademark application carries a description of the goods or services
the mark will cover. It is written to secure scope of protection, it is
dated, it is public, and it is filed by the millions of firms that never
patent. The literature that reads trademarks as innovation data has used
the filing's count, its class, its breadth and its renewal
\citep{Castaldi2020, Flikkema2019, Dinlersoz2018}; the description itself
has entered only piecemeal, as breadth counts or similarity scores. This
paper builds two research assets on the description and on the
prosecution record that surrounds it, states what each can and cannot
support, and makes both public.

The first asset is a corpus. The USPTO publishes every application since
1884 with its prosecution history, and terminal status codes are what the
literature has read from it. Status codes conflate events that the
maintenance schedule keeps apart: a registration cancelled for failing the
year-six proof of continued use and one cancelled at the year-ten renewal
share a code. Re-parsing all 13.99 million case files into 242 million
dated events assigns each cancellation to the proof it failed by its date,
and records for each filing the office actions, the responses and their
latency, counsel of record, the statutory filing basis, and the owner's
domicile. The five-year proof of continued use thereby becomes a dated,
sworn, product-level survival record available for every registered mark
more than seven years old. Owners are linked by normalized name to SEC
reporting, to Regulation~D rounds, and to PatentsView patent assignees.

The second asset is a measure. Each description is reduced to themes by a
topic model fitted once over the whole corpus, and its theme mix is
compared, by Kullback--Leibler divergence, to what its own Nice class
filed in the five years before and the five years after its date. The
average of the two divergences is the filing's \emph{atypicality}, how
unusual its language is for its industry; their difference is its
\emph{lead}, whether its language sat ahead of or behind where the
industry's language was moving. The two-sided reference is what separates
\emph{unusual} from \emph{early}, and the two components predict
different things.

The measure is validated by variation rather than by a single fit. The
identical filings are rescored under three partitions of the vocabulary
(fifty global themes, five hundred, fifty fitted per class), two
reference weightings, five window constructions, a second model seed at
fixed resolution, and with and without the 47\% of registrations whose
text is shared verbatim with another filing. Lead survives every
variation and predicts failure at the five-year proof: $+2.3$, $+2.6$ and
$+1.9$ points from the most lagging to the most leading fifth under the
three partitions, $+1.4$ points with class-by-cohort fixed effects,
drafting and filer controls and owner-clustered errors ($t = 8.3$).
Atypicality does not: its protective effect at the proof falls from $4.4$
points under global themes to $0.7$ under per-class themes, because a
global model reads a theme imported from another class as atypical, and
its sign at the SEC-reporting margin depends on how description length is
held.

Two hazards follow for any research that scores trademark text, and both
are documented on identical samples. Scoring words rather than theme
distributions measures description length: the divergence decomposes into
a surprise term and a spread term, the spread term is bounded by the
filing's distinct-term count, and term-scored atypicality correlates with
log distinct terms at $-0.68$. Firm-level correlations built on word
scoring, which are the kind the literature reports, dissolve or reverse
under theme scoring on the same firms; the signal they carry is unsigned
lexical specificity, not vocabulary position. A third caution concerns
the sample rather than the score: applicants who refile an abandoned
application rewrite toward their class's typical language from both
tails, so every registered description is in part a conformed one, and
pooling granted with abandoned filings mixes the examiner's selection
with the market's.

What the measure predicts about pioneering, about industry waves, and
about which firms reach the capital markets is the subject of a companion
paper; here those outcomes appear only as validation, and the paper makes
no claim about mechanism. Section~\ref{sec:record} describes the corpus,
the steps of the record and the three outcomes it dates.
Section~\ref{sec:measure} builds the measure and shows, on named filings,
what themes resolve and what they miss. Section~\ref{sec:related} settles
nomenclature against the innovation literature's existing terms.
Section~\ref{sec:validation} tests the measure against registration, the
five-year proof, patenting and SEC reporting. Section~\ref{sec:robust}
reports what changes under every alternative construction tried.
Section~\ref{sec:practices} states the practices that follow for
trademark-text research, and Section~\ref{sec:access} describes the
public artifacts.
